Define the \(K[x]\)-module on \(V\) inherited from a homomorphism \(T \in \text{End}(V)\),
where \(V\) is originally a vector space over \(K\). Note some essential properties.

The action of a \(f(x) \in K[x]\) on an element \(v \in V\) is given by \(f(x)v = f(T)(v)\)
The structure is also denoted \((V, T)\).

Suppose \((V, T), (W, S)\) are such \(K[x]\)-modules and \(\phi : (V, T) \to (W, S)\)
a \(K[x]\)-module hom. Then \(\phi : V \to W\) is also a vector space hom., also
\(\phi(T(v)) = \phi(x(v)) = x\phi(v) = S(\phi(v))\), since \(T = x\) in \((V, T)\).
Hence if \(\phi\) is an isomorphism it holds \(T = \phi^{-1} \circ S \phi\), which motivates similar operators.

If \((V, T)\) is free then \(\text{dim}_{K}((V)) = \infty\) as vector space and otherwise 
\((V, T)\) must be torsion and can be decomposed via the invariant/elementary factors, which are in this case
polynomials that is

\((V, T) \simeq K[x]/a_1(x) \oplus K[x]/a_2(x) \oplus \dots \oplus K[x]/a_n(x)\), with \(a_1(x) | a_2(x) | \dots | a_n(x)\).
These are unique if we enforce, that the \(a_i\) are monic. Since \(a_n(x)\) is the annhilator of \((V, T)\) the minimal polynomial 
\(m_T(x)\) is the largest invariant factor of \(V\). Additionally \((m_T(x)) = \ker{\phi}, \phi : K[x] \to K[T]\).


Derive the companion/Frobenius matrix

Suppose we want to represent \(T\) in a particular suitable basis of \(F(x)/a(x)\), where \(a(x) = x^k + b_{k-1}x^{k-1} \dots + b_0\) is the largest
invariant factor aka the minimal polynomial of \(T\). This basis is given by \(1, \bar{x}, \bar{x}^2, \dots, \bar{x}^{n-1}\)
with \(\bar{x} = x \text{mod} (a(x))\), as \(T = x\) the action of \(T\) on these elements is given by \(1 \to \bar{x}, \bar{x} \to \bar{x}^2 \dots
\bar{x}^{n-1} \to \bar{x}^k = -b_0 - b_1\bar{x} - \dots - b_{k-1}\bar{x}^{k-1}\), since \(a(\bar{x}) = 0\) in \(F[x]/(a(x))\). Thus
we may write \(T\) as:

\(
\begin{bmatrix}
0 & 0 & 0 & \dots & -b_0 \\ 
1 & 0 & 0 & \dots & -b_1 \\ 
0 & 1 & 0 & \dots & -b_2 \\ 
\vdots & \dots & \dots & \dots & \dots \\ 
0 & 0 & \dots & 1 & -b_{k-1}
\end{bmatrix}
\)


Derive the rational canonical form

Build the companion matrix \(\mathcal{C}_i\) of \(T\) for each \(a_i\) in the invariant factor decomposition
\((V, T) \simeq K[x]/a_1(x) \oplus K[x]/a_2(x) \oplus \dots \oplus K[x]/a_n(x)\)
and write them as 
\(
\begin{bmatrix}
C_1 & 0 & \dots & 0 \\ 
0 & C_2 & \dots & 0 \\ 
\vdots & \dots & \dots & \dots \\
0 & 0 & 0 & C_n \\
\end{bmatrix}
\)
in block notation.
This is the rational canonical form of \(T\).
Moreover if linear transformations \(S, T\) induce isomorphic \(K[x]\)-modules then they have the same
rational canonical form.

The special case is, if \(T \in \text{Mat}_{n}(F)\) for some field \(F\) then \(T\) is similar to a diagonal matrix \(D\)
in rational canoncial form, \(T = P D P^{-1} \).



Denote the field \(F = R/(p)\), since \((p)\) is maximal in \(R\).
It holds


Motivate the Jordan normal form

We know that \(T\) has a rational canonical form in \((V, T)\), with invariant factors \(a_1(x) | a_2(x) | \dots | a_n(x)\),
according to the elementary divisor theorem, we can instead consider the irreducible polynomial factors of these polynomials.
In the best case scenario these are linear factors, s.t. \(a_j(x) = \Pi_{i \in I}(x - \lambda_i)^{\alpha_i}\) for each \(j\).
These immediately gives a basic requirement we have for the existence of Jordan canonical form, the characteristic polynomial
(of which all these \(a_i(x)\) are divisors\) must split over \(K\).

Now consider the basis \((\bar{x} - \lambda)^{k-1}, (\bar{x} - \lambda)^{k-2}, \dots, \bar{x} - \lambda, 1\) of \(K[x]/(x - \lambda)^k\).
Then the action of \(T\) is given by 
\(x 1 \mapsto \lambda 1 + (\bar{x} - \lambda), x (\bar{x} - \lambda) \mapsto \lambda(\bar{x} - \lambda + (\bar{x} - \lambda)^2, \dots, x (\bar{x} - \lambda)^{k-1} = \lambda (\bar{x} - \lambda)^{k-1}\)
Hence we can represent \(T\) as the matrix 
\(
\begin{bmatrix}
\lambda & 1 & \dots & 0 \\ 
0 & \lambda & 1 & \dots \\ 
\vdots & \dots & \dots & \dots \\
0 & 0 & \lambda & 1 \\
0 & 0 & 0 & \lambda \\
\end{bmatrix}
\)
This is the Jordan block of size \(k\) for the eigenvalue \(\lambda\) similarly to the rational canonical form,
we can built a direct sum of the Jordan blocks for each \(K[x]/(a_i(x)\),
to obtain the Jordan canonical form (note that there is no \(1\) above a \(\lambda\) if the block is only of size \(1\))

In particular note that if \(P^{-1}AP = J\) (which always exists if \(T \in \text{Mat}_{n}(K)\)), then \(AP = PJ\).
Denote by \(p_i\) the \(i\)-th column of \(P\), if we calculate for each \(\lambda_i\) entry in \(J\), \((A - \lambda_i I_n)p_i = v\), where \(v = 0\) if \(p_i\) is regular
eigenvector or \(p_{i-1}\) if its a generalized eigenvector, hence we would have not a full basis for \(V\) if one of the \(p_i\) is a generalized
eigenvector (lin. dependent), hence the repeated \(\lambda\) entries with a \(1\) above, to migitate this issue. Moreover the eigenspaces \(T_\lambda = \{ v \in V | \exists r, (x - \lambda)^{r}v = 0 \}\),
can be seen as the subspaces of \(V\) given through the elementary divisors representation \(F[x]/(x - \lambda)^r\).


State the 4 fundamental spaces of a matrix \(A \in \text{Mat}_{n \times m}{K}\)

We have the
1. Column space \(\text{col}(A) = \text{span}(\{a_1, a_2, \dots, a_m\})\), where \(a_i\) is the \(i\)th-column of \(A\)
1. Row space, \(\text{row}(A) = \text{span}(\{a_1, a_2, \dots, a_n\})\), where \(a_i\) is the \(i\)th-row of \(A\)
1. kernel/right null space, \(\text{ker}(A) = \{ v \in K^M | Av = 0\}\)
1. left null space, \(\text{coker}(A) = \{ v \in K^n | vA = 0\}\)


We have \(\text{dim}(\text{col}(A)) = \text{dim}(\text{row}(A))\) which also implies \(\text{dim}(\text{im}(A)) = \text{dim}(\text{im}(A^{T}))\) (equivalently the rank of the matrices is the same).
The rank-nullity theorem asserts \(\text{dim}(\text{col}(A)) + \text{dim}(\text{ker}(A)) = m\) (so the dimension of the domain of \(A\)).
The rank is bounded by \(\text{rank}(A) \leq \text{min}(m, n)\).


Is there a simple way to determine the rational/Jordan canonical form for \(\text{dim} \in \{2, 3\}\)?
Moreover how can the one derive rational canonical form from the Jordan normal form and vice verca?


Yes, after finding the characteristic and minimal polynomial (the latter can be found from the prior by trial and error) of the linear transformation \(T\), 
the divisibility relations, determine already the invariant factors and with that their elementary divisors (which are the linear monomials to the highest power, which divide the invariant factor).

E.g. if \(T\) has invariant factors \(\{(x - 1)(x - 3)^{3}\}\), \(\{(x - 1)(x - 2)(x - 3)^3\}\), \(\{(x - 1)(x - 2)(x - 3)^3\}\) then elementary divisors are in turn \(\{(x - 1), (x - 3)^3\}\), \(\{(x - 1), (x - 2), (x - 3)^3\}\) etc.

Conversely if the elementary divisors are (with repetition) \(\{(x - 1)^3, (x - 1)^{3}, (x - 1)\}\), \(\{(x - 5)^2), (x - 5)\}\), \(\{(x - 4)\}\), one can arange them in a grid, where each row
lists in ascending order all the available powers of the same eigenvalue and eigenvalues with less powers fill up their first columns with \(1\)'s to achieve the same row length.
The product of the columns are then invariant factors, in this case \(\{(x - 1)\}\), \(\{(x - 1)^3(x - 5)^2\}\) \(\{(x - 1)^3(x - 5)\}\), \(\{(x - 1)^3(x - 5)^2(x - 4)\}\)


Define a \(B\)-algebra \(A\) over a commutative ring \(B\) and name some examples:

An (associative) unital \(B\)-algebra \(A\) (which is usually referred to) is a ring,
which can be seen as a \(B\)-module at the same time, with compatible multiplication, that is
\(A \times A \to A\) satisfies \((ba)a' = b(aa') = a(ba')\).

More precisely there is a ring hom. \(\phi : B \to Z(A)\), where \(Z(A)\) is the center of \(A\), s.t.
\(b \cdot a\) is defined as \(b \cdot a = \phi(b) a\).

The polynomial ring \(R[x1, \dots, x_n]\) is an \(R\)-algebra.
Each subring \(B'\) of \(B\) defines via the restriction of \(\phi' : B' \to Z(A)\) a \(B'\)-algebra
\(\text{Mat}_n(R)\) is an \(R\)-algebra which is usually not commutative.


Define and give examples of finitely generated and finite algebra's.

If there are finitely many \(a_1, \dots, a_n \in A\)  s.t. \(A = B[a_1, \dots, a_n]\) with \(B\) the
underlying ring and \(A\) the \(B\)-algebra, then \(A\) is finitely generated.

If \(A\) is finitely generated as a \(B\)-module and also a \(B\)-algebra, then \(A\) is finite.
In particular
\(A = Ba_1 + \dots + Ba_n\)

\(R\) is a finite \(R\)-algebra.
\(R[x_1, \dots, x_n]\) is a finitely gen. \(R\)-algebra, but not finite.

The inherited \(K[X]\) module from a linear transformation \(\phi\) on a finite dimensional \(K\)-vector space
is a finite algebra.


Define integral elements of a ring \(A\) over a ring \(B\).

Let \(B \subset A\) be commutative, \(a \in A\) is integral over \(B\) if there is a monic \(p(x) \in B[X]\) with 
\(p(a) = 0\). 
\(A\) is integral over \(B\) if this holds for all elements of \(A\).

Note that if \(A, B\) are fields these notions coincide with algebraic elements and normal extensions.
\(\text{End}(V)\) is integral over \((V, \phi)\) with finite dim. \(V\).


Prove and state how subalgebra's behave in regard to finiteness.

Let \(C \subset B \subset A\) be commutative rings

1. If \(B\) is a finite \(C\)-algebra and \(A\) a finite \(B\)-algebra, then \(A\) is a finite \(C\)-algebra.
2. If \(A\) is a finite \(B\)-algebra, then \(A\) is integral over \(B\).
3. If \(a \in A\) is integral over \(B\), then \(B[a]\) is a finite \(B\)-algebra. 

Proof:
(1) Let \(\{b_n\}_{n \in [m]}\) be the generating set of \(B\) as \(C\)-module and \(\{a_n\}_{n \in [m]}\) of \(A\) as \(B\)-module.
Then \(A = Ba_1 + \dots + Ba_n = (Cb_1 + \dots + Cb_n)a_1 + \dots + (Cb_1 + \dots + Cb_n)a_n\) which gives the desired result.
(2) Take any \(\alpha \in A\), use the determinant trick on \(I = B\) and \(\phi(x) = \alpha x\) (thus \(\phi(A) \subset BA\), implying
that \(\alpha\) is integral over \(B\), note that this particular \(\phi\) has powers that look like just a regular polynomial in \(x\).
with coefficients in \(B\) times a power of \(\alpha\), (setting \(x = 1\) gives then just \(\alpha\), which is valid since the specific \(p \in B[\phi]\)
is \(0\) for all arguments).
(3). Since \(\a\) is integral over \(B\) it satisfies some \(p(a) = a^{n} + b_{n-1}a^{n-1} + \dots + b_0 = 0\), hence \(B[a] = B + \dots + B a^{n-1}\).

A few direct consequences are

1. Lemma of Nakayama for algebra's:
Let \(A \neq 0\) a finite \(B\)-alebgra, then for all proper ideals \(I\) of \(B\), \(IA \neq A\).
Proof: By contraposition, if \(IA = A\) then the determinant trick gives for \(\phi = \text{id}\) that \(1 + i_{n-1} + \dots + i_0 = 0\) holds,
but then \(1 \in I\), so \(I\) is not proper.

2. Let \(A\) be a field, \(B \subset A\) a subring, s.t. \(A\) is a finite \(B\)-algebra. Then \(B\) is also a field.
From the prev. (2), \(b^{-1} \in A\) satisfies some \(p(b^{-1}) = 0\), shifting term the highest order term to the other side and 
multiplying by \(b^{n-1}\) gives that \(b^{-1} \in B\).


State and prove the existence of a coord. transformation of multinomial over a infinite field, which has a highest order
term in the last coordinate.

Let \(K\) be an infinite field \(f \in K[x_1, \dots, x_n]\), \(f \neq 0\), \(0 < d = \text{deg}f\), then there is a transformation
\(x'_i = x_i - \alpha_i x_n, \alpha_i \in K\) s.t. \(f(x'_1 + \alpha_1 x_n, \dots, x'_{n-1} + \alpha_{n-1} x_n, x_n) \in K[x_1', \dots, x_n]\)
has a \(cx^{d}_n\) term for \(0 \neq c \in K\).

Take any linear transformation of the prescribed form and note that \(f = F_d + G\), where \(F_d\) is homogenous, only with terms of \(\text{deg} = d\)
and \(G\) contains all remaining terms. Since the coordinate transformation leaves degree fixed we have 
\(f(x'_1 + \alpha_1 x_n, \dots, x'_{n-1} + \alpha_{n-1} x_n, x_n) = x^{d}_n F_d(\alpha_1, \dots, \alpha_{n-1}, 1) + \text{lower order terms or not containing x_n}\),
since \(K\) is infinite and \(F_d\) not trivial, there are \(\alpha_i \in K\), s.t. \(F_d(\alpha_1, \dots, \alpha_{n-1}, 1) \neq = 0\), thus we set \(c = F_d(\alpha_1, \dots, \alpha_{n-1}, 1)\) to obtain the desired result.

Define alebgraic ind. elements of an algebra.

Let \(A\) be a \(K\)-algebra, where \(K\) is a field and \(\{a_1, \dots, a_n\} \subset A\) the elements in question. 
If there's no non-trivial polynomial \(f \in K[t_1, \dots, t_n]\), with \(f(a_1, \dots, a_n) = 0\) then the elements are algebraically independent.

Prove and state the Noether Normalization Lemma

Let \(K\) be an infinite field \(A = K[a_1, \dots, a_n]\) a finitely generated \(K\)-algebra, then there are \(y_1, \dots y_m \in A\) with \(m \leq n\)
s.t.
1. \(y_1, \dots, y_m\) are alg. ind. over \(K\)
2. \(A\) is a finite \(K[y_1, \dots, y_m]\)-algebra.

We proceed via induction, take the generators \(a_1, \dots, a_n\) and consider the evaluation homomorphism 
\(\phi : K[x_1, \dots, x_n] \to K[a_1, \dots, a_n]\) with \(I = \text{ker}(\phi)\), if \(I = 0\) we are already done, so assume
\(I \neq 0\). If \(n = 1\) it follows \(f(a_1) = 0\), setting \(m = 0\) gives with the fact
3. If \(a \in A\) is integral over \(B\), then \(B[a]\) is a finite \(B\)-algebra
the desired result.
Let \(n > 1\), from the existence lemma of certain coord. transformations there are \(a'_i = a_i - \alpha_i a_n\)
s.t. \(f(a'_1 + \alpha_1 x_n, \dots, x_n) \in K[a'_1, \dots, x_n]\) has a \(c x_n^{d}\) term. Set 
\(A' = K[a_1', \dots, a'_{n-1}] \subset A\)
\(F(x_n) = c^{-1} f(a'_1 + \alpha_1 x_n, \dots, x_n) \in A'[x_n]\)
then \(F(x_n)\) is monic and \(F(a_n) = 0\), hence \(a_n\) is integral over \(A'\).
Applying the induction hypothesis on \(A'\) makes it a finite \(K[y_1, \dots, y_m]\)-algebra, but \(a_n\) is integral
so \(A = A'[a_n]\) is also a finite \(A'\)-algebra and the inheritance of finite algebras asserts also that \(A\) is a \(K[y_1, \dots, y_m]\)-algebra, 
yielding the desired result.


State and prove the weak Nullstellensatz (Zariski's lemma)
Let \(K\) be an infinite field, \(A\) a fin. gen. \(K\)-algebra, which is also a field,
then \(A\) is algebraic over \(K\).

According to Noether's normalization lemma, \(A\) is a finite \(B = K[y_1, \dots, y_m]\)-algebra, for some
\(y_1, \dots y_m \in A\). Moreover \(B\) is a field, since \(A\) is finite over \(B\) and due to the fact:

Let \(A\) be a field, \(B \subset A\) a subring, s.t. \(A\) is a finite \(B\)-algebra. Then \(B\) is also a field.

But then also \(m = 0\) over \(K\), hence \(A\) is finite field extension of \(K\), which must be algebraic.

